\begin{recipe}{Broccolisoep met geitenkaas}
\kicker[Voorgerecht,Vegetarisch,Frans,Kerst]{Voorgerecht · vegetarisch · Frans · kerst}
\index[register]{Broccolisoep met geitenkaas}
\index[register]{Broccoli}
\index[register]{Geitenkaas}

\meta{20 minuten \dvd Eén pan}

\lettrine{D}{eze} romige broccolisoep maakten we ooit als voorgerecht met kerst.
Het recept komt oorspronkelijk van Gordon Ramsay \cite{ramsay-broccolisoep}: met
maar een paar ingrediënten sta je zo aan tafel, en de geitenkaas en het
amandelschaafsel maken hem toch feestelijk.

\blockrule{Ingrediënten}
\begin{ingredients}
  \ing{1 kg}{broccoli, in roosjes}\index[register]{Broccoli}
  \ing{5 schijfjes}{zachte geitenkaas}\index[register]{Geitenkaas}
  \ing{800 ml}{groentebouillon}
  \ing{50 g}{amandelschaafsel}
  \ingb{olijfolie}
  \ingb{versgemalen peper en zout}
\end{ingredients}

\blockrule{Bereiding}
\tip{Gebruik een staafmixer, geen keukenmachine: die laatste maakt de soep
eerder grof dan fluweelzacht. Wil je hem helemaal glad, wrijf de soep na
het pureren nog door een fijne zeef met de achterkant van een lepel, net
als bij de tomatensoep.}
\begin{steps}
  \item Rooster het amandelschaafsel ongeveer 2 minuten op middelhoog vuur in
        een droge koekenpan tot het lichtbruin en geurig is. Zet apart.
  \item Was de broccoli en verdeel in roosjes.
  \item Breng de groentebouillon in een soeppan aan de kook en kook de
        broccoli op laag vuur 7 tot 8 minuten, tot hij zacht maar nog
        felgroen is. Giet af boven een kom, bewaar het kookvocht en doe de
        broccoli terug in de pan.
  \item Pureer de broccoli met een scheutje kookvocht met een staafmixer tot
        een gladde puree. Voeg zo nodig extra kookvocht toe voor een
        fluwelige textuur.
  \item Warm de soep op middelhoog vuur op en verdun eventueel met wat
        kookvocht. Breng op smaak met peper en zout.
  \item Schep de soep in kommen en garneer met de geitenkaas en het
        amandelschaafsel. Maal er wat peper over en besprenkel met de
        olijfolie.
\end{steps}
\end{recipe}
